\newpage
\phantomsection
\addcontentsline{toc}{section}{Список использованных источников}
\section*{СПИСОК ИСПОЛЬЗОВАННЫХ ИСТОЧНИКОВ}
\begin{enumerate}
    \item Ломакин А. А. Центробежные и осевые насосы. 2-е изд., перераб. и доп. Москва; Ленинград: Машиностроение, 1966. 364 с.
    \item Михайлов А. К., Малюшенко В. В. Лопастные насосы. Теория, расчёт и конструирование. Москва: Машиностроение, 1977. 288 с.
    \item Ломакин А. А. Расчёт критического числа оборотов и условий обеспечения динамической устойчивости роторов гидравлических машин высокого давления с учётом сил в уплотнениях // Энергомашиностроение. 1958. № 4. С. 1-8.
    \item Childs D. W. Turbomachinery Rotordynamics: Phenomena, Modeling, and Analysis. New York: Wiley, 1993.
    \item Lund J. W. Review of the Concept of Dynamic Coefficients for Fluid Film Journal Bearings // ASME Journal of Tribology. 1987. Vol. 109, no. 1. P. 37-41.
    \item Никитин Г. А. Щелевые и лабиринтные уплотнения гидроагрегатов. Москва: Машиностроение, 1982.
    \item Уплотнения и уплотнительная техника: справочник / под общ. ред. А. И. Голубева, Л. А. Кондакова. Москва: Машиностроение, 1986.
    \item Марцинковский В. А. Бесконтактные уплотнения роторных машин. Москва: Машиностроение, 1980.
    \item Марцинковский В. А. Щелевые уплотнения: теория и практика. Сумы: Изд-во СумГУ, 2005.
    \item Black H. F., Jenssen D. N. Dynamic hybrid bearing characteristics of annular controlled leakage seals // Proceedings of the Institution of Mechanical Engineers. 1969. Vol. 184. P. 92-100.
    \item Hirs G. G. A bulk-flow theory for turbulence in lubricant films // Journal of Lubrication Technology. 1973. Vol. 95, no. 2. P. 137-146.
    \item San Andrés L. Analysis of turbulent annular seals with fluid inertia effects // ASME Journal of Tribology. 1991. Vol. 113. P. 302-309.
    \item ГОСТ ИСО 1940-1-2007. Вибрация. Требования к качеству балансировки жёстких роторов. Часть 1. Определение допустимого остаточного дисбаланса. Москва: Стандартинформ, 2008.
    \item Xiong G., Zhang J., Mao Z., Wang Z., Wang H., Lian S., Jiang Z. Dynamic misalignment effects on performance of dynamically loaded journal bearings // International Journal of Mechanical Sciences. 2024. Vol. 264. Article 108839.
    \item Xie Z., Jiao J., Zhao B., Zhang J., Xu F. Theoretical and experimental research on the effect of bi-directional misalignment on the static and dynamic characteristics of a novel bearing // Mechanical Systems and Signal Processing. 2024. Vol. 208. Article 111041.
    \item Кельзон А. С., Журавлёв Ю. Н., Январёв Н. В. Расчёт и конструирование роторных машин. Ленинград: Машиностроение, 1977. 288 с.
    \item Тимошенко С. П., Янг Д. Х., Уивер У. Колебания в инженерном деле / пер. с англ. под ред. Э. И. Григолюка. Москва: Машиностроение, 1985. 472 с.
    \item Vance J. M., Zeidan F. Y., Murphy B. G. Machinery Vibration and Rotordynamics. Hoboken: John Wiley \& Sons, 2010.
    \item Genta G. Dynamics of Rotating Systems. New York: Springer, 2005. 658 p.
    \item Adams M. L. Rotating Machinery Vibration: From Analysis to Troubleshooting. 2nd ed. Boca Raton: CRC Press, 2010.
    \item API Standard 610. Centrifugal Pumps for Petroleum, Petrochemical and Natural Gas Industries. 12th ed. Washington: American Petroleum Institute, 2021; ISO 13709:2009. Centrifugal pumps for petroleum, petrochemical and natural gas industries. Geneva: ISO, 2009.
    \item ГОСТ Р ИСО 20816-3-2023. Вибрация. Измерение и оценка вибрации машин. Часть 3. Машины промышленного назначения номинальной мощностью более 15 кВт и номинальными скоростями от 120 об/мин до 30000 об/мин. Москва: Стандартинформ, 2023.
    \item API Recommended Practice 684. API Standard Paragraphs Rotordynamic Tutorial: Lateral Critical Speeds, Unbalance Response, Stability, Train Torsionals, and Rotor Balancing. 4th ed. Washington: American Petroleum Institute, 2014. Reaffirmed 2022.
    \item ГОСТ 801-78. Сталь подшипниковая. Технические условия. Москва: Издательство стандартов, 1978.
    \item Mota J. A., Maldonado D. J. G., Valério J. V., Ritto T. G. Complete modeling of hydrodynamic bearings with a boundary parameterization approach // Mechanism and Machine Theory. 2021. Vol. 162. Article 104350.
    \item Степанов А. И. Центробежные и осевые насосы. Москва: Машгиз, 1960.
    \item Gu C., Meng X., Wang S., Ding X. Modeling a hydrodynamic bearing with provision for misalignments and textures // Journal of Tribology. 2020. Vol. 142. Article 041801.
    \item ТУ 3631-003-00217389-96. Насосы центробежные секционные типа ЦНС. Технические условия.
    \item Колесо рабочее 1 ступени ЦНС 90.00.00.012. Чертёж детали, УГНТУ МП-98-02.
    \item ГОСТ 31381-2009. Насосы динамические. Методы испытаний. Москва: Стандартинформ, 2010.
\end{enumerate}
